% Options for packages loaded elsewhere
% Options for packages loaded elsewhere
\PassOptionsToPackage{unicode}{hyperref}
\PassOptionsToPackage{hyphens}{url}
%
\documentclass[
  ignorenonframetext,
  aspectratio=169,
  t]{beamer}
\newif\ifbibliography
\usepackage{pgfpages}
\setbeamertemplate{caption}[numbered]
\setbeamertemplate{caption label separator}{: }
\setbeamercolor{caption name}{fg=normal text.fg}
\beamertemplatenavigationsymbolsempty
% remove section numbering
\setbeamertemplate{part page}{
  \centering
  \begin{beamercolorbox}[sep=16pt,center]{part title}
    \usebeamerfont{part title}\insertpart\par
  \end{beamercolorbox}
}
\setbeamertemplate{section page}{
  \centering
  \begin{beamercolorbox}[sep=12pt,center]{section title}
    \usebeamerfont{section title}\insertsection\par
  \end{beamercolorbox}
}
\setbeamertemplate{subsection page}{
  \centering
  \begin{beamercolorbox}[sep=8pt,center]{subsection title}
    \usebeamerfont{subsection title}\insertsubsection\par
  \end{beamercolorbox}
}
% Prevent slide breaks in the middle of a paragraph
\widowpenalties 1 10000
\raggedbottom
\AtBeginPart{
  \frame{\partpage}
}
\AtBeginSection{
  \ifbibliography
  \else
    \frame{\sectionpage}
  \fi
}
\AtBeginSubsection{
  \frame{\subsectionpage}
}
\usepackage{iftex}
\ifPDFTeX
  \usepackage[T1]{fontenc}
  \usepackage[utf8]{inputenc}
  \usepackage{textcomp} % provide euro and other symbols
\else % if luatex or xetex
  \usepackage{unicode-math} % this also loads fontspec
  \defaultfontfeatures{Scale=MatchLowercase}
  \defaultfontfeatures[\rmfamily]{Ligatures=TeX,Scale=1}
\fi
\usepackage{lmodern}

\ifPDFTeX\else
  % xetex/luatex font selection
\fi
% Use upquote if available, for straight quotes in verbatim environments
\IfFileExists{upquote.sty}{\usepackage{upquote}}{}
\IfFileExists{microtype.sty}{% use microtype if available
  \usepackage[]{microtype}
  \UseMicrotypeSet[protrusion]{basicmath} % disable protrusion for tt fonts
}{}
\makeatletter
\@ifundefined{KOMAClassName}{% if non-KOMA class
  \IfFileExists{parskip.sty}{%
    \usepackage{parskip}
  }{% else
    \setlength{\parindent}{0pt}
    \setlength{\parskip}{6pt plus 2pt minus 1pt}}
}{% if KOMA class
  \KOMAoptions{parskip=half}}
\makeatother

\usepackage{color}
\usepackage{fancyvrb}
\newcommand{\VerbBar}{|}
\newcommand{\VERB}{\Verb[commandchars=\\\{\}]}
\DefineVerbatimEnvironment{Highlighting}{Verbatim}{commandchars=\\\{\}}
% Add ',fontsize=\small' for more characters per line
\usepackage{framed}
\definecolor{shadecolor}{RGB}{241,243,245}
\newenvironment{Shaded}{\begin{snugshade}}{\end{snugshade}}
\newcommand{\AlertTok}[1]{\textcolor[rgb]{0.68,0.00,0.00}{#1}}
\newcommand{\AnnotationTok}[1]{\textcolor[rgb]{0.37,0.37,0.37}{#1}}
\newcommand{\AttributeTok}[1]{\textcolor[rgb]{0.40,0.45,0.13}{#1}}
\newcommand{\BaseNTok}[1]{\textcolor[rgb]{0.68,0.00,0.00}{#1}}
\newcommand{\BuiltInTok}[1]{\textcolor[rgb]{0.00,0.23,0.31}{#1}}
\newcommand{\CharTok}[1]{\textcolor[rgb]{0.13,0.47,0.30}{#1}}
\newcommand{\CommentTok}[1]{\textcolor[rgb]{0.37,0.37,0.37}{#1}}
\newcommand{\CommentVarTok}[1]{\textcolor[rgb]{0.37,0.37,0.37}{\textit{#1}}}
\newcommand{\ConstantTok}[1]{\textcolor[rgb]{0.56,0.35,0.01}{#1}}
\newcommand{\ControlFlowTok}[1]{\textcolor[rgb]{0.00,0.23,0.31}{\textbf{#1}}}
\newcommand{\DataTypeTok}[1]{\textcolor[rgb]{0.68,0.00,0.00}{#1}}
\newcommand{\DecValTok}[1]{\textcolor[rgb]{0.68,0.00,0.00}{#1}}
\newcommand{\DocumentationTok}[1]{\textcolor[rgb]{0.37,0.37,0.37}{\textit{#1}}}
\newcommand{\ErrorTok}[1]{\textcolor[rgb]{0.68,0.00,0.00}{#1}}
\newcommand{\ExtensionTok}[1]{\textcolor[rgb]{0.00,0.23,0.31}{#1}}
\newcommand{\FloatTok}[1]{\textcolor[rgb]{0.68,0.00,0.00}{#1}}
\newcommand{\FunctionTok}[1]{\textcolor[rgb]{0.28,0.35,0.67}{#1}}
\newcommand{\ImportTok}[1]{\textcolor[rgb]{0.00,0.46,0.62}{#1}}
\newcommand{\InformationTok}[1]{\textcolor[rgb]{0.37,0.37,0.37}{#1}}
\newcommand{\KeywordTok}[1]{\textcolor[rgb]{0.00,0.23,0.31}{\textbf{#1}}}
\newcommand{\NormalTok}[1]{\textcolor[rgb]{0.00,0.23,0.31}{#1}}
\newcommand{\OperatorTok}[1]{\textcolor[rgb]{0.37,0.37,0.37}{#1}}
\newcommand{\OtherTok}[1]{\textcolor[rgb]{0.00,0.23,0.31}{#1}}
\newcommand{\PreprocessorTok}[1]{\textcolor[rgb]{0.68,0.00,0.00}{#1}}
\newcommand{\RegionMarkerTok}[1]{\textcolor[rgb]{0.00,0.23,0.31}{#1}}
\newcommand{\SpecialCharTok}[1]{\textcolor[rgb]{0.37,0.37,0.37}{#1}}
\newcommand{\SpecialStringTok}[1]{\textcolor[rgb]{0.13,0.47,0.30}{#1}}
\newcommand{\StringTok}[1]{\textcolor[rgb]{0.13,0.47,0.30}{#1}}
\newcommand{\VariableTok}[1]{\textcolor[rgb]{0.07,0.07,0.07}{#1}}
\newcommand{\VerbatimStringTok}[1]{\textcolor[rgb]{0.13,0.47,0.30}{#1}}
\newcommand{\WarningTok}[1]{\textcolor[rgb]{0.37,0.37,0.37}{\textit{#1}}}

\usepackage{longtable,booktabs,array}
\newcounter{none} % for unnumbered tables
\usepackage{calc} % for calculating minipage widths
\usepackage{caption}
% Make caption package work with longtable
\makeatletter
\def\fnum@table{\tablename~\thetable}
\makeatother
\usepackage{graphicx}
\makeatletter
\newsavebox\pandoc@box
\newcommand*\pandocbounded[1]{% scales image to fit in text height/width
  \sbox\pandoc@box{#1}%
  \Gscale@div\@tempa{\textheight}{\dimexpr\ht\pandoc@box+\dp\pandoc@box\relax}%
  \Gscale@div\@tempb{\linewidth}{\wd\pandoc@box}%
  \ifdim\@tempb\p@<\@tempa\p@\let\@tempa\@tempb\fi% select the smaller of both
  \ifdim\@tempa\p@<\p@\scalebox{\@tempa}{\usebox\pandoc@box}%
  \else\usebox{\pandoc@box}%
  \fi%
}
% Set default figure placement to htbp
\def\fps@figure{htbp}
\makeatother





\setlength{\emergencystretch}{3em} % prevent overfull lines

\providecommand{\tightlist}{%
  \setlength{\itemsep}{0pt}\setlength{\parskip}{0pt}}



 


\setbeamertemplate{frametitle continuation}{}
\setbeamertemplate{itemize items}[circle]
\usepackage{caption}
\captionsetup{labelformat=empty, justification=centering}
\usepackage{textpos}
\setlength{\TPHorizModule}{1cm}
\setlength{\TPVertModule}{1cm}
\setbeamertemplate{footline}{%
  \hfill\usebeamercolor[fg]{page number in head/foot}%
  \usebeamerfont{page number in head/foot}%
  \insertframenumber{} / \inserttotalframenumber\kern1em\vskip2pt%
}
\usepackage{xcolor}
\let\oldtexttt\texttt
\renewcommand{\texttt}[1]{\colorbox{gray!20}{\oldtexttt{#1}}}

\AtBeginSection[]{ % for centering the section slides since others are top aligned
  \begin{frame}[c]
  \centering
  \Large\insertsectionhead
  \end{frame}
}
\makeatletter
\@ifpackageloaded{caption}{}{\usepackage{caption}}
\AtBeginDocument{%
\ifdefined\contentsname
  \renewcommand*\contentsname{Table of contents}
\else
  \newcommand\contentsname{Table of contents}
\fi
\ifdefined\listfigurename
  \renewcommand*\listfigurename{List of Figures}
\else
  \newcommand\listfigurename{List of Figures}
\fi
\ifdefined\listtablename
  \renewcommand*\listtablename{List of Tables}
\else
  \newcommand\listtablename{List of Tables}
\fi
\ifdefined\figurename
  \renewcommand*\figurename{Figure}
\else
  \newcommand\figurename{Figure}
\fi
\ifdefined\tablename
  \renewcommand*\tablename{Table}
\else
  \newcommand\tablename{Table}
\fi
}
\@ifpackageloaded{float}{}{\usepackage{float}}
\floatstyle{ruled}
\@ifundefined{c@chapter}{\newfloat{codelisting}{h}{lop}}{\newfloat{codelisting}{h}{lop}[chapter]}
\floatname{codelisting}{Listing}
\newcommand*\listoflistings{\listof{codelisting}{List of Listings}}
\makeatother
\makeatletter
\makeatother
\makeatletter
\@ifpackageloaded{caption}{}{\usepackage{caption}}
\@ifpackageloaded{subcaption}{}{\usepackage{subcaption}}
\makeatother

\usepackage{bookmark}
\IfFileExists{xurl.sty}{\usepackage{xurl}}{} % add URL line breaks if available
\urlstyle{same}
% fallback for those not using the hyperref driver hyperxmp:
\makeatletter
\@ifundefined{xmpquote}{\newcommand{\xmpquote}[1]{#1}}{}
\makeatother
\hypersetup{
  pdftitle={Introduction to Basic and Advanced Statistical Modelling},
  pdfauthor={Alexandre Courtiol},
  hidelinks,
  pdfcreator={LaTeX via pandoc}}


\title{\texorpdfstring{Introduction to Basic and Advanced Statistical
Modelling}{Introduction to Basic and Advanced Statistical Modelling}}
\subtitle{\texorpdfstring{Using R (solutions to
exercises)}{Using R (solutions to exercises)}}
\author{Alexandre Courtiol}
\date{}

\begin{document}
\frame{\titlepage}

\renewcommand*\contentsname{Table of contents}
\begin{frame}[allowframebreaks]
  \frametitle{Table of contents}
  \setcounter{tocdepth}{2}
  \tableofcontents
\end{frame}
\setcounter{tocdepth}{2}
\tableofcontents
}

\section{Reading data}\label{reading-data}

\begin{frame}[fragile]{Practice 4}
\protect\phantomsection\label{practice-4}
\footnotesize

\begin{Shaded}
\begin{Highlighting}[]
\NormalTok{territories }\OtherTok{\textless{}{-}} \FunctionTok{read\_csv}\NormalTok{(}\StringTok{"data\_MHB/data{-}territories.csv"}\NormalTok{)}
\end{Highlighting}
\end{Shaded}

\begin{verbatim}
Rows: 1152306 Columns: 10
-- Column specification ------------------------------------------------------------------
Delimiter: ","
chr (2): nameeg, namelt
dbl (8): species_id, euring_id, samplearea_id, year, territories_visit_1, territories_...

i Use `spec()` to retrieve the full column specification for this data.
i Specify the column types or set `show_col_types = FALSE` to quiet this message.
\end{verbatim}

\normalsize
\end{frame}

\section{Wrangling data}\label{wrangling-data}

\begin{frame}[fragile]{Practice 5a}
\protect\phantomsection\label{practice-5a}
Use \texttt{select()} and/or \texttt{filter()} and/or \texttt{slice()}
and/or \texttt{arrange()} to determine:

\begin{enumerate}
\tightlist
\item
  the largest number of recorded species across all years
\end{enumerate}

\footnotesize

\begin{Shaded}
\begin{Highlighting}[]
\NormalTok{species\_count }\SpecialCharTok{|\textgreater{}} 
    \FunctionTok{arrange}\NormalTok{(}\FunctionTok{desc}\NormalTok{(n\_recorded\_species)) }\SpecialCharTok{|\textgreater{}} 
    \FunctionTok{slice}\NormalTok{(}\DecValTok{1}\NormalTok{) }\SpecialCharTok{|\textgreater{}} 
    \FunctionTok{select}\NormalTok{(n\_recorded\_species)}
\end{Highlighting}
\end{Shaded}

\normalsize

or

\footnotesize

\begin{Shaded}
\begin{Highlighting}[]
\NormalTok{species\_count }\SpecialCharTok{|\textgreater{}} 
    \FunctionTok{filter}\NormalTok{(n\_recorded\_species }\SpecialCharTok{==} \FunctionTok{max}\NormalTok{(n\_recorded\_species)) }\SpecialCharTok{|\textgreater{}} 
    \FunctionTok{select}\NormalTok{(n\_recorded\_species)}
\end{Highlighting}
\end{Shaded}

\normalsize
\end{frame}

\begin{frame}[fragile]{Practice 5b}
\protect\phantomsection\label{practice-5b}
Use \texttt{select()}, \texttt{filter()}, \texttt{slice()} and/or
\texttt{arrange()} to determine:

\begin{enumerate}
\setcounter{enumi}{1}
\tightlist
\item
  the largest number of recorded species in 2025
\end{enumerate}

\footnotesize

\begin{Shaded}
\begin{Highlighting}[]
\NormalTok{species\_count }\SpecialCharTok{|\textgreater{}} 
    \FunctionTok{filter}\NormalTok{(year }\SpecialCharTok{==} \DecValTok{2025}\NormalTok{) }\SpecialCharTok{|\textgreater{}} 
    \FunctionTok{filter}\NormalTok{(n\_recorded\_species }\SpecialCharTok{==} \FunctionTok{max}\NormalTok{(n\_recorded\_species)) }\SpecialCharTok{|\textgreater{}} 
    \FunctionTok{select}\NormalTok{(n\_recorded\_species)}
\end{Highlighting}
\end{Shaded}

\normalsize

but NOT this:

\footnotesize

\begin{Shaded}
\begin{Highlighting}[]
\NormalTok{species\_count }\SpecialCharTok{|\textgreater{}} 
    \FunctionTok{filter}\NormalTok{(year }\SpecialCharTok{==} \DecValTok{2025}\NormalTok{, n\_recorded\_species }\SpecialCharTok{==} \FunctionTok{max}\NormalTok{(n\_recorded\_species)) }\SpecialCharTok{|\textgreater{}}
    \FunctionTok{select}\NormalTok{(n\_recorded\_species)}
\end{Highlighting}
\end{Shaded}

\normalsize

Do you understand why?
\end{frame}

\begin{frame}[fragile]{Practice 6a}
\protect\phantomsection\label{practice-6a}
Use the functions we have seen to obtain:

\begin{enumerate}
\tightlist
\item
  the range of median elevation according to the number of visits
\end{enumerate}

\footnotesize

\begin{Shaded}
\begin{Highlighting}[]
\NormalTok{species\_count }\SpecialCharTok{|\textgreater{}} 
    \FunctionTok{summarise}\NormalTok{(}\AttributeTok{median\_elevation\_min =} \FunctionTok{min}\NormalTok{(median\_elevation),}
              \AttributeTok{median\_elevation\_max =} \FunctionTok{max}\NormalTok{(median\_elevation),}
              \AttributeTok{.by =}\NormalTok{ n\_visits)}
\end{Highlighting}
\end{Shaded}

\begin{verbatim}
# A tibble: 2 x 3
  n_visits median_elevation_min median_elevation_max
     <dbl>                <dbl>                <dbl>
1        3                  250                 2150
2        2                 1450                 2750
\end{verbatim}

\normalsize
\end{frame}

\begin{frame}[fragile]{Practice 6b}
\protect\phantomsection\label{practice-6b}
Use the functions we have seen to obtain:

\begin{enumerate}
\setcounter{enumi}{1}
\tightlist
\item
  the 3 most recorded species for each year
\end{enumerate}

\footnotesize

\begin{Shaded}
\begin{Highlighting}[]
\NormalTok{territories }\SpecialCharTok{|\textgreater{}} 
    \FunctionTok{summarise}\NormalTok{(}\AttributeTok{territories\_sum =} \FunctionTok{sum}\NormalTok{(territories\_total),}
              \AttributeTok{.by =} \FunctionTok{c}\NormalTok{(year, nameeg, namelt)) }\SpecialCharTok{|\textgreater{}} \CommentTok{\# aggregate across samplearea\_id}
    \FunctionTok{slice\_max}\NormalTok{(territories\_sum, }\AttributeTok{n =} \DecValTok{3}\NormalTok{, }\AttributeTok{by =}\NormalTok{ year) }\SpecialCharTok{|\textgreater{}} \CommentTok{\# pick top 3 per year}
    \FunctionTok{select}\NormalTok{(year, nameeg, namelt, territories\_sum)}
\end{Highlighting}
\end{Shaded}

\begin{verbatim}
# A tibble: 81 x 4
   year nameeg             namelt            territories_sum
  <dbl> <chr>              <chr>                       <dbl>
1  1999 Eurasian Chaffinch Fringilla coelebs            5775
2  1999 Common Blackbird   Turdus merula                3132
3  1999 Coal Tit           Periparus ater               3108
4  2000 Eurasian Chaffinch Fringilla coelebs            6695
# i 77 more rows
\end{verbatim}

\normalsize
\end{frame}

\begin{frame}[fragile]{Practice 6c (1)}
\protect\phantomsection\label{practice-6c-1}
Use the functions we have seen to obtain:

\begin{enumerate}
\setcounter{enumi}{2}
\tightlist
\item
  the number of observations for \emph{Sturnus vulgaris} for each year
\end{enumerate}

\footnotesize

\begin{Shaded}
\begin{Highlighting}[]
\NormalTok{territories }\SpecialCharTok{|\textgreater{}} 
    \FunctionTok{filter}\NormalTok{(namelt }\SpecialCharTok{==} \StringTok{"Sturnus vulgaris"}\NormalTok{) }\SpecialCharTok{|\textgreater{}} 
    \FunctionTok{summarise}\NormalTok{(}\AttributeTok{territories\_sum =} \FunctionTok{sum}\NormalTok{(territories\_total, }\AttributeTok{na.rm =} \ConstantTok{TRUE}\NormalTok{), }\CommentTok{\# NAs present}
              \AttributeTok{.by =}\NormalTok{ year)}
\end{Highlighting}
\end{Shaded}

\begin{verbatim}
# A tibble: 27 x 2
   year territories_sum
  <dbl>           <dbl>
1  1999             691
2  2000             920
3  2001             911
4  2002             887
# i 23 more rows
\end{verbatim}

\normalsize
\end{frame}

\begin{frame}[fragile]{Practice 6c (2)}
\protect\phantomsection\label{practice-6c-2}
Use the functions we have seen to obtain:

\begin{enumerate}
\setcounter{enumi}{2}
\tightlist
\item
  the number of observations for \emph{Sturnus vulgaris} for each year
\end{enumerate}

\footnotesize

\begin{Shaded}
\begin{Highlighting}[]
\NormalTok{territories }\SpecialCharTok{|\textgreater{}} 
    \FunctionTok{filter}\NormalTok{(namelt }\SpecialCharTok{==} \StringTok{"Sturnus vulgaris"}\NormalTok{) }\SpecialCharTok{|\textgreater{}} 
    \FunctionTok{drop\_na}\NormalTok{(territories\_total) }\SpecialCharTok{|\textgreater{}} \CommentTok{\# alternative not to use na.rm below}
    \FunctionTok{summarise}\NormalTok{(}\AttributeTok{territories\_sum =} \FunctionTok{sum}\NormalTok{(territories\_total), }\CommentTok{\# NAs present}
              \AttributeTok{.by =}\NormalTok{ year)}
\end{Highlighting}
\end{Shaded}

\begin{verbatim}
# A tibble: 27 x 2
   year territories_sum
  <dbl>           <dbl>
1  1999             691
2  2000             920
3  2001             911
4  2002             887
# i 23 more rows
\end{verbatim}

\normalsize
\end{frame}

\begin{frame}[fragile]{Practice 7}
\protect\phantomsection\label{practice-7}
Try to use \texttt{pivot\_longer()} to get back to the original long
format of the data from the output of the following code:

\footnotesize

\begin{Shaded}
\begin{Highlighting}[]
\NormalTok{territories }\SpecialCharTok{|\textgreater{}} 
    \FunctionTok{filter}\NormalTok{(year }\SpecialCharTok{==} \DecValTok{2025}\NormalTok{) }\SpecialCharTok{|\textgreater{}} 
    \FunctionTok{select}\NormalTok{(samplearea\_id, nameeg, territories\_total) }\SpecialCharTok{|\textgreater{}} 
    \FunctionTok{pivot\_wider}\NormalTok{(}\AttributeTok{names\_from =}\NormalTok{ nameeg, }
                \AttributeTok{values\_from =}\NormalTok{ territories\_total) }\SpecialCharTok{|\textgreater{}} 
    \FunctionTok{pivot\_longer}\NormalTok{(}\SpecialCharTok{{-}}\NormalTok{samplearea\_id,}
                 \AttributeTok{names\_to =} \StringTok{"nameeg"}\NormalTok{,}
                 \AttributeTok{values\_to =} \StringTok{"territories\_total"}\NormalTok{)}
\end{Highlighting}
\end{Shaded}

\begin{verbatim}
# A tibble: 42,282 x 3
  samplearea_id nameeg              territories_total
          <dbl> <chr>                           <dbl>
1           268 Little Grebe                        0
2           268 Great Crested Grebe                 0
3           268 Grey Heron                          0
4           268 Little Bittern                      0
# i 42,278 more rows
\end{verbatim}

\normalsize
\end{frame}

\section{Wrangling spatial data}\label{wrangling-spatial-data}

\begin{frame}[fragile]{Practice 8a}
\protect\phantomsection\label{practice-8a}
Use the functions we have seen to obtain:

\begin{enumerate}
\tightlist
\item
  the sampled area (sensu convex hull) across years
\end{enumerate}

\footnotesize

\begin{Shaded}
\begin{Highlighting}[]
\NormalTok{species\_count\_sf }\SpecialCharTok{|\textgreater{}} 
    \FunctionTok{summarise}\NormalTok{(}\AttributeTok{geometry =} \FunctionTok{st\_combine}\NormalTok{(geometry), }\AttributeTok{.by =}\NormalTok{ year) }\SpecialCharTok{|\textgreater{}} 
    \FunctionTok{mutate}\NormalTok{(}\AttributeTok{hull =} \FunctionTok{st\_convex\_hull}\NormalTok{(geometry),}
           \AttributeTok{area =} \FunctionTok{st\_area}\NormalTok{(hull),}
           \AttributeTok{area\_km2 =}\NormalTok{ units}\SpecialCharTok{::}\FunctionTok{set\_units}\NormalTok{(area, }\StringTok{"km\^{}2"}\NormalTok{)) }\SpecialCharTok{|\textgreater{}} 
    \FunctionTok{select}\NormalTok{(}\SpecialCharTok{{-}}\NormalTok{hull) }\SpecialCharTok{|\textgreater{}} 
    \FunctionTok{st\_drop\_geometry}\NormalTok{()}
\end{Highlighting}
\end{Shaded}

\begin{verbatim}
# A tibble: 27 x 3
   year         area area_km2
* <dbl>        [m^2]   [km^2]
1  1999 48899915383.   48900.
2  2000 48995663682.   48996.
3  2001 48229824468.   48230.
4  2002 48995663682.   48996.
# i 23 more rows
\end{verbatim}

\normalsize
\end{frame}

\begin{frame}[fragile]{Practice 8b (1)}
\protect\phantomsection\label{practice-8b-1}
Use the functions we have seen to obtain:

\begin{enumerate}
\setcounter{enumi}{1}
\tightlist
\item
  the IDs of the sampled areas which are most peripheral
\end{enumerate}

Simple sorting by distance to centroid (but lacks threshold)

\footnotesize

\begin{Shaded}
\begin{Highlighting}[]
\NormalTok{species\_count\_sf }\SpecialCharTok{|\textgreater{}} 
  \FunctionTok{distinct}\NormalTok{(samplearea\_id, geometry) }\SpecialCharTok{|\textgreater{}} 
  \FunctionTok{mutate}\NormalTok{(}\AttributeTok{centroid =} \FunctionTok{st\_centroid}\NormalTok{(}\FunctionTok{st\_combine}\NormalTok{(geometry)),}
         \AttributeTok{dist =} \FunctionTok{st\_distance}\NormalTok{(centroid, geometry, }\AttributeTok{by\_element =} \ConstantTok{TRUE}\NormalTok{)) }\SpecialCharTok{|\textgreater{}} 
  \FunctionTok{arrange}\NormalTok{(}\FunctionTok{desc}\NormalTok{(dist)) }\SpecialCharTok{|\textgreater{}} 
  \FunctionTok{as\_tibble}\NormalTok{() }\CommentTok{\# (for slide only {-}\textgreater{} skip metadata clutter)}
\end{Highlighting}
\end{Shaded}

\begin{verbatim}
# A tibble: 267 x 4
  samplearea_id            geometry            centroid    dist
          <dbl>         <POINT [°]>         <POINT [°]>     [m]
1           269 (6.110458 46.21013) (8.234991 46.79076) 174955.
2           398 (10.40901 46.61099) (8.234991 46.79076) 166982.
3           271 (6.184874 46.35493) (8.234991 46.79076) 164026.
4           397 (10.34234 46.82876) (8.234991 46.79076) 160429.
# i 263 more rows
\end{verbatim}

\normalsize
\end{frame}

\begin{frame}[fragile]{Practice 8b (2)}
\protect\phantomsection\label{practice-8b-2}
More useful, but more complex: intersection with concave hull outline

\footnotesize

\begin{Shaded}
\begin{Highlighting}[]
\NormalTok{species\_count\_sf }\SpecialCharTok{|\textgreater{}} 
  \FunctionTok{distinct}\NormalTok{(samplearea\_id, }\AttributeTok{.keep\_all =} \ConstantTok{TRUE}\NormalTok{) }\SpecialCharTok{|\textgreater{}} 
  \FunctionTok{mutate}\NormalTok{(}\AttributeTok{all\_points =} \FunctionTok{st\_combine}\NormalTok{(geometry), }\CommentTok{\# combine points into new geometry}
         \AttributeTok{hull =} \FunctionTok{st\_concave\_hull}\NormalTok{(all\_points, }\AttributeTok{ratio =} \FloatTok{0.2}\NormalTok{), }\CommentTok{\# create concave hull}
         \AttributeTok{perimeter =} \FunctionTok{st\_boundary}\NormalTok{(hull), }\CommentTok{\# extract outer line of the hull}
         \AttributeTok{perimeter =} \FunctionTok{st\_buffer}\NormalTok{(perimeter, }\AttributeTok{dist =} \DecValTok{1000}\NormalTok{), }\CommentTok{\# create buffer zone around it}
         \AttributeTok{on\_perim =} \FunctionTok{st\_intersects}\NormalTok{(perimeter, geometry, }\CommentTok{\# test if points fall into buffer}
                                  \AttributeTok{by\_element =} \ConstantTok{TRUE}\NormalTok{)) }\OtherTok{{-}\textgreater{}}\NormalTok{ data\_perimeter\_sf }
  
\NormalTok{data\_perimeter\_sf }\SpecialCharTok{|\textgreater{}} 
    \FunctionTok{filter}\NormalTok{(on\_perim) }\SpecialCharTok{|\textgreater{}} 
    \FunctionTok{pull}\NormalTok{(samplearea\_id) }\CommentTok{\# pull() extracts a column}
\end{Highlighting}
\end{Shaded}

\begin{verbatim}
 [1] 269 271 273 277 278 279 282 283 286 289 291 294 302 307 308 311 323 325 340 352 359
[22] 369 370 376 384 385 388 392 393 394 395 396 397 398 400 408 411 418 427 434 439 456
[43] 460 461 462 469 475 476 478 502 503 515 522 525 531
\end{verbatim}

\normalsize
\end{frame}

\begin{frame}[fragile]{Practice 8b (2)}
\protect\phantomsection\label{practice-8b-2-1}
More useful, but more complex: intersection with concave hull outline

\footnotesize

\begin{Shaded}
\begin{Highlighting}[]
\FunctionTok{theme\_set}\NormalTok{(}\FunctionTok{theme\_bw}\NormalTok{(}\AttributeTok{base\_size =} \DecValTok{20}\NormalTok{)) }\CommentTok{\# setting the theme for all plots}
\FunctionTok{ggplot}\NormalTok{(data\_perimeter\_sf) }\SpecialCharTok{+}
  \FunctionTok{geom\_sf}\NormalTok{(}\FunctionTok{aes}\NormalTok{(}\AttributeTok{geometry =}\NormalTok{ perimeter)) }\SpecialCharTok{+}
  \FunctionTok{geom\_sf}\NormalTok{(}\FunctionTok{aes}\NormalTok{(}\AttributeTok{colour =}\NormalTok{ on\_perim)) }\CommentTok{\# geometry = geometry is implicit}
\end{Highlighting}
\end{Shaded}

\begin{center}
\includegraphics[width=\linewidth,height=0.5\textheight,keepaspectratio]{Rintro_solutions_files/figure-beamer/unnamed-chunk-14-1.pdf}
\end{center}

\normalsize
\end{frame}

\section{Plotting data}\label{plotting-data}

\begin{frame}[fragile]{Practice 10 (1)}
\protect\phantomsection\label{practice-10-1}
\begin{itemize}
\tightlist
\item
  combine data wrangling and plotting to illustrate the change in
  species abundance across years for the top 6 most encountered species
\end{itemize}

\footnotesize

\begin{Shaded}
\begin{Highlighting}[]
\NormalTok{territories }\SpecialCharTok{|\textgreater{}} 
    \FunctionTok{summarise}\NormalTok{(}\AttributeTok{territories\_sum =} \FunctionTok{sum}\NormalTok{(territories\_total, }\AttributeTok{na.rm =} \ConstantTok{TRUE}\NormalTok{),}
              \AttributeTok{.by =} \FunctionTok{c}\NormalTok{(nameeg, namelt)) }\SpecialCharTok{|\textgreater{}} 
    \FunctionTok{slice\_max}\NormalTok{(territories\_sum, }\AttributeTok{n =} \DecValTok{6}\NormalTok{) }\SpecialCharTok{|\textgreater{}} 
    \FunctionTok{pull}\NormalTok{(namelt) }\OtherTok{{-}\textgreater{}}\NormalTok{ top6sp}

\NormalTok{top6sp}
\end{Highlighting}
\end{Shaded}

\begin{verbatim}
[1] "Fringilla coelebs"  "Sylvia atricapilla" "Turdus merula"      "Periparus ater"    
[5] "Erithacus rubecula" "Passer domesticus" 
\end{verbatim}

\normalsize
\end{frame}

\begin{frame}[fragile]{Practice 10 (2)}
\protect\phantomsection\label{practice-10-2}
\begin{itemize}
\tightlist
\item
  combine data wrangling and plotting to illustrate the change in
  species abundance across years for the top 6 most encountered species
\end{itemize}

\footnotesize

\begin{Shaded}
\begin{Highlighting}[]
\NormalTok{territories }\SpecialCharTok{|\textgreater{}} 
    \FunctionTok{filter}\NormalTok{(namelt }\SpecialCharTok{\%in\%}\NormalTok{ top6sp) }\SpecialCharTok{|\textgreater{}} 
    \FunctionTok{summarise}\NormalTok{(}\AttributeTok{territories\_sum =} \FunctionTok{sum}\NormalTok{(territories\_total, }\AttributeTok{na.rm =} \ConstantTok{TRUE}\NormalTok{),}
              \AttributeTok{.by =} \FunctionTok{c}\NormalTok{(nameeg, namelt, year)) }\OtherTok{{-}\textgreater{}}\NormalTok{ top6\_across\_time}

\NormalTok{top6\_across\_time}
\end{Highlighting}
\end{Shaded}

\begin{verbatim}
# A tibble: 162 x 4
  nameeg   namelt          year territories_sum
  <chr>    <chr>          <dbl>           <dbl>
1 Coal Tit Periparus ater  1999            3108
2 Coal Tit Periparus ater  2000            3568
3 Coal Tit Periparus ater  2001            2803
4 Coal Tit Periparus ater  2002            2959
# i 158 more rows
\end{verbatim}

\normalsize
\end{frame}

\begin{frame}[fragile]{Practice 10 (3a)}
\protect\phantomsection\label{practice-10-3a}
\begin{itemize}
\tightlist
\item
  combine data wrangling and plotting to illustrate the change in
  species abundance across years for the top 6 most encountered species
\end{itemize}

\footnotesize

\begin{Shaded}
\begin{Highlighting}[]
\FunctionTok{ggplot}\NormalTok{(top6\_across\_time) }\SpecialCharTok{+}
  \FunctionTok{geom\_line}\NormalTok{(}\FunctionTok{aes}\NormalTok{(}\AttributeTok{y =}\NormalTok{ territories\_sum, }\AttributeTok{x =}\NormalTok{ year, }\AttributeTok{colour =}\NormalTok{ nameeg))}
\end{Highlighting}
\end{Shaded}

\begin{center}
\includegraphics[width=0.85\linewidth,height=\textheight,keepaspectratio]{Rintro_solutions_files/figure-beamer/unnamed-chunk-17-1.pdf}
\end{center}

\normalsize
\end{frame}

\begin{frame}[fragile]{Practice 10 (3b)}
\protect\phantomsection\label{practice-10-3b}
\begin{itemize}
\tightlist
\item
  combine data wrangling and plotting to illustrate the change in
  species abundance across years for the top 6 most encountered species
\end{itemize}

\footnotesize

\begin{Shaded}
\begin{Highlighting}[]
\FunctionTok{ggplot}\NormalTok{(top6\_across\_time) }\SpecialCharTok{+}
  \FunctionTok{geom\_line}\NormalTok{(}\FunctionTok{aes}\NormalTok{(}\AttributeTok{y =}\NormalTok{ territories\_sum, }\AttributeTok{x =}\NormalTok{ year)) }\SpecialCharTok{+}
  \FunctionTok{facet\_wrap}\NormalTok{(}\SpecialCharTok{\textasciitilde{}}\NormalTok{ nameeg, }\AttributeTok{scales =} \StringTok{"free\_y"}\NormalTok{)}
\end{Highlighting}
\end{Shaded}

\begin{center}
\includegraphics[width=\linewidth,height=0.6\textheight,keepaspectratio]{Rintro_solutions_files/figure-beamer/unnamed-chunk-18-1.pdf}
\end{center}

\normalsize
\end{frame}




\end{document}
